\chapter{Research Methodology}\label{ch:method}

This chapter describes how the research was conducted. It sets out the philosophical stance and overall design, the software development process, the requirements that the system was built to satisfy, the procedures for collecting drone data and for the preceding mock-data phase, the experimental design used to compare controllers, the metrics and statistical tests, the threats to validity, the tools and project plan, and the ethical, legal, social and professional issues raised by the work.

% ------------------------------------------------------------------
\section{Research Philosophy and Approach}\label{sec:method-philosophy}

The research adopted a \emph{positivist} philosophy. Its central claims (for example, that one controller produces lower mean waiting time than another at a given junction) are empirical propositions about an external system that can be measured objectively and tested against data. Accordingly, the approach was \emph{quantitative} and \emph{deductive}: hypotheses derived from traffic-control theory, notably that pressure-based control stabilises queues more effectively than demand-blind timing~\cite{varaiya2013maxpressure}, were operationalised as measurable differences in waiting time, queue length and throughput, and tested statistically.

Within this stance the work combined two established strategies. The first is \emph{design science}, in which knowledge is generated by building and evaluating purposeful artefacts; here the artefacts were \dtx{}, \vls{} and \prism{}, and their utility was judged against explicit requirements. The second is \emph{experimental simulation}: because intervening on a live public junction was neither safe nor permissible for an undergraduate project, controlled experiments were performed on a calibrated microscopic model of the junction~\cite{dowling2004toolbox}. Simulation permits what field trials cannot: identical demand can be replayed under each controller using common random numbers, so that observed differences are attributable to the controller alone; sensor faults can be injected at chosen intensities; and many replications can be run cheaply. The price of this choice is that conclusions are only as valid as the model, which is why a substantial part of the methodology was devoted to obtaining real observations and calibrating against them.

An interpretivist or qualitative approach, such as interviewing traffic police officers about their manual control decisions, was considered and would have offered valuable insight into local practice. It was not adopted as a primary method because the research questions concern measurable performance, but informal observation of manual control informed the problem framing in \cref{ch:intro}.

% ------------------------------------------------------------------
\section{Research Design Overview}\label{sec:method-design}

The research was organised into eight phases, shown in \cref{fig:method-phases}. The literature review and requirements analysis fixed the problem and the measurable criteria. Phase~1 then developed and verified the complete simulation, sensing and control stack on synthetic mock demand. Phase~2 collected real drone video, after which \dtx{} extracted trajectories and calibration quantities, the SUMO model was calibrated, and the final experiments and analysis were performed. Calibration was itself iterative, with model parameters re-tuned until the acceptance criteria of \cref{sec:method-metrics} were met. The dashed arrow indicates the principal feedback loop of the design: any anomaly observed in the experiments was traced back to controller development, corrected and re-verified on mock data before being re-run on real demand.

\begin{figure}[H]
\centering
\resizebox{\textwidth}{!}{%
\begin{tikzpicture}[
  node distance=0.55cm and 0.6cm,
  box/.style={draw, rounded corners=3pt, align=center, minimum width=3.3cm, minimum height=1.0cm, font=\small, fill=#1},
  box/.default={blue!8},
  arr/.style={-{Stealth[length=2.5mm]}, thick},
  fb/.style={-{Stealth[length=2.5mm]}, dashed, gray}
]
\node[box] (lit) {Literature\\review};
\node[box, right=of lit] (req) {Requirements\\analysis};
\node[box=orange!12, right=of req] (p1) {Phase 1: mock-data\\development \&\\verification};
\node[box=green!12, below=of p1] (p2) {Phase 2: drone\\data collection\\(2 Oct 2026)};
\node[box=green!12, left=of p2] (dtx) {\dtx{} trajectory\\extraction};
\node[box=green!12, left=of dtx] (cal) {SUMO model\\calibration};
\node[box=purple!10, below=of cal] (exp) {Controller\\experiments};
\node[box=purple!10, right=of exp] (ana) {Statistical analysis\\\& sensitivity};
\node[box=gray!12, right=of ana] (wr) {Reporting};
\draw[arr] (lit) -- (req);
\draw[arr] (req) -- (p1);
\draw[arr] (p1) -- (p2);
\draw[arr] (p2) -- (dtx);
\draw[arr] (dtx) -- (cal);
\draw[arr] (cal) -- (exp);
\draw[arr] (exp) -- (ana);
\draw[arr] (ana) -- (wr);
\draw[fb] (exp.south) -- ++(0,-0.45) -| ($(wr.east)+(0.45,0)$) |- node[pos=0.25, right, font=\scriptsize, align=left]{controller fixes,\\re-verified on\\mock data} (p1.east);
\end{tikzpicture}}
\caption{Phases of the research design. Mock-data development (Phase~1) preceded drone data collection and calibration (Phase~2).}
\label{fig:method-phases}
\end{figure}

% ------------------------------------------------------------------
\section{Development Methodology}\label{sec:method-dev}

Software was developed using an \emph{iterative and incremental} methodology organised as two-week, agile-style sprints. A plan-driven waterfall model was rejected because several key uncertainties (detector accuracy on small motorcycles, the realism of sublane filtering, and the stability of the controller under noisy sensing) could only be resolved by building and testing; committing to a full design up front would have risked late and expensive rework. A pure exploratory approach was equally unsuitable, because the drone flight date was a hard external milestone.

Each sprint had a single increment as its goal (for example, ``\vls{} with noise and dropout'', or ``\prism{} gap-out and starvation override''), a backlog maintained as GitHub issues, and an end-of-sprint review with the supervisor. Increments were integrated into a single Git repository whose structure mirrors the system architecture (\texttt{extraction/}, \texttt{sumo/}, \texttt{sensing/}, \texttt{control/}, \texttt{experiments/} and \texttt{tests/}). Development followed a test-first discipline for safety-critical logic: tests asserting that no two conflicting approaches could ever be green simultaneously, that amber and all-red intervals were always inserted, and that no approach could be starved beyond $T_{\text{starve}}$ were written before the corresponding controller code, and the full \texttt{pytest} suite was run before every merge to the main branch. Every experiment configuration was stored in version-controlled YAML so that any figure in \cref{ch:eval} can be traced to a specific commit.

% ------------------------------------------------------------------
\section{Requirements Analysis}\label{sec:method-req}

Requirements were elicited from three sources: the literature on adaptive control and simulation practice~\cite{koonce2008timing,dowling2004toolbox}, the physical and operational characteristics of the Maviththara junction, and the research questions. They are listed in \cref{tab:fr,tab:nfr}; each is traced to the design in \cref{ch:design} and verified in \cref{ch:eval}.

\begin{table}[H]
\centering
\caption{Functional requirements.}
\label{tab:fr}
\small
\begin{tabularx}{\textwidth}{@{}lXl@{}}
\toprule
\textbf{ID} & \textbf{Requirement} & \textbf{Component}\\
\midrule
FR1 & Stabilise drone video to a reference frame to remove hover drift. & \dtx{}\\
FR2 & Detect and classify seven vehicle classes (motorcycle, three-wheeler, car, van, bus, lorry, bicycle) in aerial frames. & \dtx{}\\
FR3 & Track vehicles across frames and convert positions to metric ground coordinates. & \dtx{}\\
FR4 & Derive per-class flows, turning ratios, saturation flow, queue lengths and per-vehicle waiting times, and export SUMO route files. & \dtx{}\\
FR5 & Represent the junction in SUMO with 2+2 lanes per arm, left-hand traffic and sublane filtering of two- and three-wheelers. & Model\\
FR6 & Simulate a 360\textdegree{} roadside LiDAR with range limit, noise, dropout, occlusion, size-class confusion and clutter. & \vls{}\\
FR7 & Estimate per-approach queued, arriving and downstream PCU from \vls{} output only. & \prism{}\\
FR8 & Select phases and green durations every 1~s using the pressure-based law, subject to $G_{\min}$, $G_{\max}$, amber, all-red and starvation constraints. & \prism{}\\
FR9 & Fall back to the fixed-time plan when sensor data are unavailable. & \prism{}\\
FR10 & Run batch experiments for all controllers, scenarios and seeds, and compute all evaluation metrics and statistical tests. & Harness\\
\bottomrule
\end{tabularx}
\end{table}

\begin{table}[H]
\centering
\caption{Non-functional requirements.}
\label{tab:nfr}
\small
\begin{tabularx}{\textwidth}{@{}lXl@{}}
\toprule
\textbf{ID} & \textbf{Requirement} & \textbf{Criterion}\\
\midrule
NFR1 & \emph{Safety}: conflicting approaches are never simultaneously green; every change includes 3~s amber and 2~s all-red. & Zero violations in tests\\
NFR2 & \emph{Real-time performance}: a control decision completes well within the 1~s control interval. & $<$ 10~ms per decision\\
NFR3 & \emph{Fairness}: no approach waits longer than the starvation cap. & Max red $\le$ 120~s\\
NFR4 & \emph{Explainability}: every decision is computable by hand from logged state; no learned model in the control loop. & Design review\\
NFR5 & \emph{Reproducibility}: all runs are deterministic given a seed and configuration under version control. & Bit-identical re-runs\\
NFR6 & \emph{Privacy}: no personally identifiable information is stored at any stage. & Data audit\\
\bottomrule
\end{tabularx}
\end{table}

% ------------------------------------------------------------------
\section{Data Collection Methodology}\label{sec:method-data}

\subsection{Choice of instrument}
Several data-collection methods were considered for observing demand and performance at the junction. Manual roadside counts are inexpensive but cannot record trajectories, waiting times or queue lengths for hundreds of vehicles simultaneously, and are error-prone for dense motorcycle streams. Roadside cameras observe only a partial, oblique view in which vehicles occlude each other. Inductive loops are absent at Maviththara and in any case measure only point presence. An unmanned aerial vehicle (UAV) hovering above the junction, by contrast, provides a single nadir view of all four approaches and the junction box simultaneously, with minimal occlusion and a geometry that maps to the ground plane by a single homography. Drone video is now an established source of ground-truth naturalistic trajectories, as demonstrated by the highD~\cite{krajewski2018highd} and pNEUMA~\cite{barmpounakis2020pneuma} datasets, and it doubles as an independent ground truth against which the simulated LiDAR's view can be compared.

\subsection{Flight procedure}
Two flights were performed on Thursday, 2~October~2026, a normal working and school day: Session~M from 07:00 to 08:00, covering the morning peak with severe congestion, and Session~A from 13:30 to 14:30, representing moderate demand. A DJI-class quadcopter was flown to an altitude of approximately \est{90~m} and held in a stationary hover with the gimbal pointing vertically downwards (nadir), recording at \est{3840$\times$2160} pixels and \est{30~fps}. At this altitude the field of view covered roughly the final \est{120~m} of each approach, and the ground sample distance was approximately \est{2.6~cm/px}. Because a single battery does not last an hour, each session was recorded in \est{three} segments with brief battery changes, during which the take-off point was retained so that every segment could be stabilised to the same reference frame; the short gaps (\est{under 90~s} each) were treated as censoring boundaries in \dtx{}. Ground control points (lane-marking corners and stop-line ends) were identified for georeferencing, with metric coordinates derived from measured lane widths of 3.5~m and satellite imagery. Weather on both sessions was dry.

\subsection{Regulatory compliance and anonymity}
Flights were planned in accordance with the Civil Aviation Authority of Sri Lanka (CAASL) implementing standards for remotely piloted aircraft systems~\cite{caasl_drone}, including registration of the aircraft, observation of the permitted altitude ceiling, maintenance of visual line of sight, avoidance of controlled airspace and of flight directly over gatherings of people, and notification of local authorities; confirmation of the specific permissions obtained is recorded in the appendices. Anonymity was ensured by design: at \est{90~m} altitude, faces and number plates occupy only a few pixels and are unidentifiable, and the \dtx{} pipeline retained only anonymous numerical trajectories (track identifier, class, time, position and speed). Raw video was stored on an encrypted drive, was not published, and is scheduled for deletion after examination.

% ------------------------------------------------------------------
\section{Mock-Data Phase}\label{sec:method-mock}

Phase~1 used synthetic demand so that development could proceed for several months before real data were available. Demand was split evenly across the four approaches at three levels: low (300~veh/h per approach), medium (500~veh/h) and high (700~veh/h). Every approach used turning ratios of 25\% left, 50\% through and 25\% right, and the vehicle mix approximated the national registration profile~\cite{dmt_stats}: 40\% motorcycles, 20\% three-wheelers, 25\% cars, 7\% vans, 5\% buses and 3\% lorries.

The symmetric design was intentional. With identical demand on all approaches, a correct controller should allocate approximately equal total green to each approach over an hour; marked asymmetry therefore revealed bugs in state estimation or map-matching. The three demand levels exercised distinct regimes: at 300~veh/h the gap-out logic dominates, at 500~veh/h pressure switching is most active, and at 700~veh/h the junction approaches saturation under split phasing, stressing $G_{\max}$, starvation and spillback logic. Phase~1 served three purposes: (i) \emph{verification}, through unit and integration tests of phase safety, \vls{} noise statistics, pressure calculation, gap-out and starvation behaviour; (ii) \emph{parameter tuning}, through a grid search over $\theta$, $\gamma$, $\delta$, $\kappa$ and $G_{\max}$ (seeds 1--5); and (iii) a \emph{controlled benchmark}, the medium-demand scenario being retained in the final results as a reference point free of calibration error.

% ------------------------------------------------------------------
\section{Simulation-Based Experimental Design}\label{sec:method-experiments}

\subsection{Scenarios and controllers}
Three demand scenarios were evaluated: \textbf{Mock} (500~veh/h per approach), \textbf{Afternoon} (moderate demand from Session~A, \est{2,280~veh/h} in total) and \textbf{Morning} (heavy demand from Session~M, \est{3,420~veh/h}). Real-demand scenarios used the time-varying 5-minute flows and turning ratios exported by \dtx{}. In each scenario five controllers were compared, all using the same four-phase split structure and the same 3~s amber and 2~s all-red intervals:
\begin{itemize}[itemsep=1pt]
  \item \textbf{B1 Fixed-time}: the observed current plan (cycle \est{120~s}).
  \item \textbf{B2 Actuated}: SUMO's built-in gap-based actuated controller~\cite{sumo_actuated} with minimum 10~s and maximum 60~s greens, representing conventional detector-based actuation.
  \item \textbf{B3 Webster}: a fixed-time plan re-optimised with Webster's method~\cite{webster1958} from the \dtx{} demand, representing the best achievable \emph{demand-blind} plan and ensuring that \prism{} is not credited merely for correcting a poorly tuned existing plan.
  \item \textbf{\prism{}}: the proposed controller observing the junction through the noisy \vls{}.
  \item \textbf{\prism{}-GT}: the same controller fed ground-truth vehicle states, an oracle that isolates the cost of imperfect sensing.
\end{itemize}

\subsection{Replication and tuning protocol}
Each run consisted of a 300~s warm-up period, discarded so that the junction was populated before measurement, followed by 3600~s of measurement. Vehicles were included in metrics if they entered the network during the measurement window. Stochasticity in SUMO (vehicle insertion, driver imperfection $\sigma$ and speed factors) was controlled by the random seed. To avoid optimistic bias from tuning and evaluating on the same realisations, seeds were partitioned: seeds~1--5 were used exclusively for parameter tuning, and seeds~6--15 (ten seeds) were held out for final evaluation. Within each evaluation seed, all five controllers experienced identical demand realisations (common random numbers), which makes the comparisons paired and greatly reduces the variance of the estimated differences. In total the final evaluation comprised $3 \times 5 \times 10 = 150$ simulation runs, plus the sensitivity analyses.

\subsection{Sensitivity analyses}
Two families of one-factor-at-a-time sensitivity experiments were run on the Afternoon scenario. Sensor sensitivity varied the base position noise $\sigma_0 \in \{0, 0.1, 0.3, 0.5, 1.0\}$~m and the dropout probability $p \in \{0, 0.05, 0.1, 0.2, 0.3\}$. Parameter sensitivity varied the hysteresis $\theta \in \{0, 0.1, 0.25, 0.5\}$ and $G_{\max} \in \{45, 60, 75, 90\}$~s. All other parameters were held at their design values (\cref{app:params}).

% ------------------------------------------------------------------
\section{Evaluation Metrics}\label{sec:method-metrics}

Let $\mathcal{V}$ denote the set of $N$ vehicles that completed their trips during the measurement window, and let vehicle $i$ have speed $v_i(t_k)$ at simulation step $t_k$ with step length $\Delta t = 0.5$~s.

\paragraph{Average waiting time (primary metric).} A vehicle is waiting at a step when its speed is below $v_{\text{w}} = 0.1$~m/s, the threshold used by SUMO and, by design, by \dtx{}:
\begin{equation}
  w_i = \sum_{k} \Delta t \cdot \mathbb{1}\!\left[v_i(t_k) < v_{\text{w}}\right], \qquad
  \bar{W} = \frac{1}{N}\sum_{i \in \mathcal{V}} w_i .
  \label{eq:wait}
\end{equation}

\paragraph{95th-percentile waiting time.} $W_{95}$ is the 95th percentile of $\{w_i\}$, capturing the experience of the worst-served vehicles and hence fairness.

\paragraph{Average travel time.} With entry and exit times $t_i^{\text{in}}$ and $t_i^{\text{out}}$ on the approximately 400~m arms, $\bar{T} = \frac{1}{N}\sum_i (t_i^{\text{out}} - t_i^{\text{in}})$.

\paragraph{Average queue length.} For approach $a$ the queue length $L_a(t)$ is the distance from the stop line to the rear of the last halting vehicle (speed $< v_{\text{w}}$) in a contiguous queue. The mean over the set of approaches $\mathcal{A}$ and measurement horizon $T_m$ is
\begin{equation}
  \bar{L} = \frac{1}{|\mathcal{A}|\,T_m}\sum_{a \in \mathcal{A}} \int_{0}^{T_m} L_a(t)\,\mathrm{d}t ,
  \label{eq:queue}
\end{equation}
evaluated in discrete time; the maximum queue $L_{\max} = \max_{a,t} L_a(t)$ is also reported.

\paragraph{Throughput.} $\Theta$ is the number of vehicles completing their trips per hour of measurement. Under demand below capacity throughput equals demand, so differences are expected mainly in the Morning scenario.

\paragraph{Stops per vehicle.} $\bar{S} = \frac{1}{N}\sum_i s_i$, where $s_i$ counts transitions from moving to waiting (SUMO's \texttt{waitingCount}).

\paragraph{CO$_2$ emissions.} Total CO$_2$ per hour, $E = \sum_i \sum_k e_{\text{CO}_2}(v_i(t_k), a_i(t_k), c_i)\,\Delta t$, where $e_{\text{CO}_2}$ is the HBEFA emission rate~\cite{hbefa41} for vehicle class $c_i$ at the instantaneous speed and acceleration.

\paragraph{Relative improvement.} For any metric $X$, the change of controller $C$ relative to fixed-time is
\begin{equation}
  \Delta_X(C) = \frac{X_C - X_{\text{B1}}}{X_{\text{B1}}} \times 100\% .
  \label{eq:relchange}
\end{equation}

\paragraph{Statistical testing.} Because only ten paired observations are available per comparison and normality of the differences cannot be assumed, the non-parametric Wilcoxon signed-rank test~\cite{wilcoxon1945} was used. For each evaluation seed $j$ the paired difference $d_j = \bar{W}_{\text{\prism{}},j} - \bar{W}_{\text{B},j}$ was computed; the absolute differences were ranked and the statistic
\begin{equation}
  T^{+} = \sum_{j:\, d_j > 0} \operatorname{rank}\left(|d_j|\right)
  \label{eq:wilcoxon}
\end{equation}
was compared against its exact null distribution. With $n = 10$, the smallest attainable two-sided $p$-value is $2/2^{10} \approx 0.002$, so significance at the 0.01 level is attainable only when \prism{} wins on (almost) every seed. Results are reported as mean $\pm$ standard deviation over the ten seeds, together with $p$-values; a significance level of $\alpha = 0.05$ was used throughout.

\paragraph{Calibration metrics.} Simulated and observed hourly counts $M$ and $C$ were compared with the GEH statistic,
\begin{equation}
  \mathrm{GEH} = \sqrt{\frac{2(M - C)^2}{M + C}} ,
  \label{eq:geh}
\end{equation}
with the accepted criterion GEH $< 5$ for at least 85\% of movements~\cite{dowling2004toolbox}. Waiting-time fidelity was measured by the mean absolute percentage error between simulated and observed average waiting time.

% ------------------------------------------------------------------
\section{Validity and Reliability}\label{sec:method-validity}

\textbf{Internal validity} concerns whether observed differences are caused by the controllers. It was protected by paired common-random-number seeds, identical phase structures and intergreen times for all controllers, warm-up removal, and the separation of tuning and evaluation seeds. The inclusion of an optimised fixed-time baseline (B3) guards against attributing to adaptivity what is merely re-timing.

\textbf{Construct validity} concerns whether the metrics measure what is intended. Defining waiting time identically in \dtx{} and in SUMO ($v < 0.1$~m/s) ensured that the calibrated model and the observations measure the same construct. Reporting $W_{95}$, maximum red time and per-class waiting alongside the mean prevents an efficiency gain from masking unfairness.

\textbf{External validity} is the principal limitation. Results derive from one junction, observed on one weekday, in dry weather; they may not transfer to junctions with turn pockets, different phasing or substantially different composition, and the simulation omits police manual control, pedestrians and the behaviour of drivers who might adapt to a new controller. The real LiDAR is represented by a model whose error parameters, although grounded in sensor specifications~\cite{ouster_os1} and the roadside LiDAR literature~\cite{zhao2019lidar}, have not been validated against a physical unit at this site; the sensitivity analysis partly mitigates this by showing performance across a range of error levels.

\textbf{Conclusion validity} was addressed by using a non-parametric test appropriate for small paired samples and by reporting dispersion as well as means. \textbf{Reliability} was ensured by fixing software versions, seeding all random generators (SUMO, \vls{} and NumPy), storing all configurations in Git, and automating the complete pipeline from route files to tables so that every reported result can be regenerated.

% ------------------------------------------------------------------
\section{Tools and Technologies}\label{sec:method-tools}

\Cref{tab:tools} summarises the tools used and the reasons for their selection; implementation details are given in \cref{ch:impl}.

\begin{table}[H]
\centering
\caption{Tools and technologies.}
\label{tab:tools}
\small
\begin{tabularx}{\textwidth}{@{}lX@{}}
\toprule
\textbf{Tool} & \textbf{Purpose and justification}\\
\midrule
Python 3.11 & Single language for extraction, sensing, control and analysis; mature scientific ecosystem.\\
Ultralytics YOLOv8~\cite{jocher2023yolov8} & Anchor-free single-stage detector; fine-tuning from VisDrone weights~\cite{zhu2022visdrone}.\\
SAHI~\cite{akyon2022sahi} & Sliced inference for small aerial objects (motorcycles $\approx$ 20~px).\\
supervision / ByteTrack~\cite{zhang2022bytetrack} & Multi-object tracking without a re-identification network.\\
OpenCV, NumPy, SciPy, Pandas & Stabilisation (ORB, RANSAC), homography, Savitzky--Golay smoothing, data handling.\\
CVAT / Roboflow & Manual bounding-box annotation of \est{450} training frames.\\
Google Colab (T4 GPU) & Detector fine-tuning without local GPU hardware.\\
SUMO \est{1.20.0}~\cite{lopez2018sumo} & Open-source microscopic simulator with sublane model and HBEFA emissions.\\
TraCI / libsumo~\cite{wegener2008traci} & Online control of signals and retrieval of vehicle states; libsumo for batch speed.\\
OpenStreetMap~\cite{haklay2008osm}, netedit & Base network geometry and manual cleaning.\\
pytest & Unit and integration tests for safety-critical controller logic.\\
Matplotlib, Jupyter & Analysis and figures.\\
Git / GitHub & Version control, issue tracking, sprint backlog.\\
\LaTeX{} & Dissertation preparation.\\
\bottomrule
\end{tabularx}
\end{table}

% ------------------------------------------------------------------
\section{Project Management}\label{sec:method-pm}

The project ran from June to October 2026. \Cref{tab:gantt} shows the plan at half-month resolution. The critical external constraint was the drone flight, scheduled for early October; the mock-data phase was planned so that all components would be verified before that date, leaving the final weeks for extraction, calibration and experiments. Principal risks and mitigations were: (i) flight cancellation due to rain or permission delays, mitigated by retaining the mock-data results as a fallback evaluation and reserving alternative dates; (ii) poor detector accuracy on motorcycles, mitigated by sliced inference and local fine-tuning; and (iii) calibration failure, mitigated by early preparation of the calibration scripts on mock data.

\begin{table}[H]
\centering
\caption{Project timeline (each column is half a month; $\blacksquare$ indicates planned activity).}
\label{tab:gantt}
\small
\setlength{\tabcolsep}{4pt}
\begin{tabular}{@{}lcccccccccc@{}}
\toprule
& \multicolumn{2}{c}{Jun} & \multicolumn{2}{c}{Jul} & \multicolumn{2}{c}{Aug} & \multicolumn{2}{c}{Sep} & \multicolumn{2}{c}{Oct}\\
\cmidrule(lr){2-3}\cmidrule(lr){4-5}\cmidrule(lr){6-7}\cmidrule(lr){8-9}\cmidrule(lr){10-11}
\textbf{Activity} & 1 & 2 & 1 & 2 & 1 & 2 & 1 & 2 & 1 & 2\\
\midrule
Literature review & $\blacksquare$ & $\blacksquare$ & $\blacksquare$ & & & & & & & \\
Requirements \& design & & $\blacksquare$ & $\blacksquare$ & & & & & & & \\
SUMO network (mock demand) & & & $\blacksquare$ & $\blacksquare$ & & & & & & \\
\vls{} development & & & & $\blacksquare$ & $\blacksquare$ & & & & & \\
\prism{} \& baselines & & & & & $\blacksquare$ & $\blacksquare$ & $\blacksquare$ & & & \\
Mock-data verification \& tuning & & & & & & $\blacksquare$ & $\blacksquare$ & $\blacksquare$ & & \\
\dtx{} development (sample clips) & & & & & & $\blacksquare$ & $\blacksquare$ & $\blacksquare$ & & \\
Drone flights (2 Oct) & & & & & & & & & $\blacksquare$ & \\
Extraction \& calibration & & & & & & & & & $\blacksquare$ & \\
Experiments \& analysis & & & & & & & & & $\blacksquare$ & $\blacksquare$\\
Report writing & & & & & & & $\blacksquare$ & $\blacksquare$ & $\blacksquare$ & $\blacksquare$\\
\bottomrule
\end{tabular}
\end{table}

% ------------------------------------------------------------------
\section{Ethical, Legal, Social and Professional Issues}\label{sec:method-elsp}

\textbf{Ethical.} The only human-related data were drone recordings of public road users. Data minimisation was applied at source: the altitude was chosen so that individuals and number plates were unidentifiable, and only anonymous trajectories were retained. No participant was contacted and no personal data were processed, so informed consent was not required, but institutional ethical approval for the data collection was sought through the supervisor before the flights. The design choice to rely on LiDAR rather than cameras for a future deployment is itself an ethical commitment to privacy-preserving infrastructure.

\textbf{Legal.} Drone operation complied with the CAASL standards for remotely piloted aircraft~\cite{caasl_drone}. The handling of recordings was guided by the principles of Sri Lanka's Personal Data Protection Act No.~9 of 2022, notably purpose limitation, data minimisation, storage limitation and security, even though anonymised trajectories fall outside its definition of personal data. All software used is open source or used under academic licences, and OpenStreetMap data are used under the Open Database Licence with attribution.

\textbf{Social.} Adaptive control redistributes delay; a controller that reduces average waiting can, if poorly designed, do so by persistently disadvantaging a minor approach or a vehicle class. The starvation cap, the reporting of per-class and 95th-percentile waiting, and the explicit PCU weighting of buses (which carry many passengers) address this concern. The research also recognises that traffic police perform an important role at the junction; \prism{} is presented as a tool to reduce the need for manual intervention, with police override envisaged as a necessary feature of any deployment.

\textbf{Professional.} The work was conducted in accordance with the BCS Code of Conduct: in the \emph{public interest}, by prioritising safety constraints (NFR1) over efficiency and by stating the limitations of simulation-based evidence honestly; with \emph{professional competence and integrity}, by acknowledging all sources, reporting negative and sensitivity results as well as favourable ones, and not overstating the readiness of the system for deployment; and with \emph{duty to the relevant authority}, by respecting aviation regulations and institutional procedures. All code and configurations are retained under version control so that the results can be independently checked.
